%% mortal -- mortal holes of the two-word problem with self events:
%% where they come from (facts of phi) and an example, anchored comparisons,
%% the separator word and its reader, which admits self events, the coarse
%% depth bound, positive channels and their derivatives with the witness in
%% the tower, the anchored compression theorem, short mortal holes, and
%% decidability of the two-word problem under (FH). Proofs and constants are
%% in appendix-mortal.tex.

\section{Mortal Holes}\label{sec:mo}

Let $I=(A,\Bad,B,\Rej,\Guard,\GuardS,A',B',\Head)$ be an instance of the
two-word hole problem with self events (\cref{def:ex-ss2self}), with the
notation of \cref{def:ex-ss2self,def:ex-events}. A \emph{mortal hole} of $I$
is a hole $(u,t)$ whose spine ends in a mortal state and whose target ends
fewer than $N_A$ letters after its first rejected cut
(\cref{def:nm-mortal-hole}). By \cref{thm:nm-three-cases}, an instance with
forced target heads has a hole iff it has a short hole or a mortal hole. This
section bounds the length of a shortest mortal hole and thereby decides the
hole problem under (FH) (\cref{thm:mo-decide}).

\paragraph{Where mortal holes come from}
In an extracted instance $I_{pq}$, mortality depends on $\varphi$ alone. The
spine automaton is the lower automaton of $q$ in $\varphi$, and its bad states
are those that meet $\Fup\varphi$, so a spine $u$ is safe iff
$(q,u)\notin\mu_\varphi$ (\cref{tab:ex-instance}): the derived atom is not one
of the atoms that every solution of $\varphi$ contains. Facts come only from
extreme literals (\cref{tab:lf-clauses}). Hence, if $\varphi$ has no extreme
literal, then $\mu_\varphi$, $\Fup\varphi$ and $\Bad$ are empty, and for
$n\ge1$ no state is mortal, so $I_{pq}$ has no mortal hole
(\cref{rem:nm-nonblocking}). Mortal holes arise only when $\varphi$ forces
atoms. Since every safe run from a mortal state dies within $N_A$ letters
(\cref{lem:nm-thin}), a safe spine $u$ ends in a mortal state iff
$(q,uv)\in\mu_\varphi$ for every word $v$ of length $N_A$: the derived atom is
not forced, but all its long extensions are. A mortal hole therefore admits no
arbitrarily long safe continuation. The argument of \cref{sec:nm}, which replaces the continuation of a
hole by a new one taken from the infinite safe continuations of the spine,
does not apply, and the pair has to be shortened as it stands.

\begin{example}[A mortal hole]
Let $n=2$ with $f$ covariant, let $X=\{x,y,a,b,c\}$ and $Z=\varnothing$, and
\[
 \varphi=\{x\le f(x,b),\ b\le f(a,a),\ a=\Bot,\ f(y,c)\le y\},\qquad
 \psi=\{x\le f(y,y)\}.
\]
Both systems are satisfiable. Condition \CondI\ holds, since $\psi$ has no
facts, and so does \CondU: the only atom that is root-dead in $\psi$ is its
forbidden root $(x^0,\eps)$, which is also a forbidden root of $\varphi$ and
hence dead in $\varphi$. Consider $I_{y^1x^1}$.
\begin{itemize}[leftmargin=*]
\item The only fact of $\varphi$ is $(a^1,\eps)$, and $\Fup\varphi=\{a^1\}$.
  The spine automaton starts at $\downc\varphi{x^1}=\{x^1\}$, loops there on
  $0$, moves to $\{b^1\}$ on $1$, and from $\{b^1\}$ to the bad state
  $\{a^1\}$ on either letter; so $\{b^1\}$ is mortal. In terms of trees,
  $x\le f(x,b)$, $b\le f(a,a)$ and $a=\Bot$ force a leaf $\Bot$ of $x$ at a
  prefix of every word $0^k1i$.
\item The $\varphi$-component of the target automaton is $\{y^1\}$ along
  $0^*$, $\{c^1\}$ after the first letter $1$, and empty afterwards. The atoms
  $(y^1,0^k)$ are dead by support, because $(y^1,\eps)$ is a forbidden root of
  $\varphi$ and $(y^1,0^k)$ pops to it, while no atom $(y^1,w)$ is root-dead
  in $\psi$. So $J(t)$ is defined iff $t$ contains the letter $1$, and it is
  then the cut after the first $1$.
\item The guard and the self guard never apply, because the sets that they
  intersect are disjoint.
\item The pop head is empty: $\psi$ derives $(x^1,0\sigma)$ from
  $(y^1,\sigma)$ by a single push, and the head events are the pairs $(1,0)$
  with $\wsuf u1=t$.
\end{itemize}
Hence the holes of $I_{y^1x^1}$ are exactly the pairs $(0^{j+1}1,0^j1)$ with
$j\ge0$, all of them mortal, with $J(t)=\len t$; the spine and the common
suffix $0^j1$ are unbounded. These pairs are also the only pre-holes, and their
spines end in the mortal state, so the instance has holes, but only in case~(3)
of \cref{thm:nm-three-cases}. In the corresponding failure of \CondII, $\psi$
derives $(x^1,0^{j+1}1)$ from the live seed $(y^1,0^j1)$, but the least model
of $\varphi$ that contains the seed does not contain $(x^1,0^{j+1}1)$. For
$j=0$, the solution $a=c=\Bot$, $b=f(\Bot,\Bot)$, $x=f(f(\Bot,b),b)$,
$y=f(\Top,\Bot)$ of $\varphi$ is not a solution of $\psi$: $x\le f(y,y)$ would
require $f(\Bot,b)\le f(\Top,\Bot)$, hence $b\le\Bot$.
\end{example}

In a mortal hole the target ends shortly after its first rejected cut, but the
spine and the common suffix of a head event can be arbitrarily long, as in the
example. They cannot be shortened independently: events relate cuts of the
spine to cuts of the target, self events relate two cuts of the target, and a
head event demands that two suffixes be \emph{equal}. We proceed in four steps.
\begin{enumerate}[leftmargin=*]
\item We pack the pair into the \emph{separator word} $u\sep t$ over
  $\Sigma\sqcup\{\sep\}$. Events and self events become comparisons of this
  word, and a head event becomes an \emph{anchored} comparison, one whose
  matched factor is followed by the separator. Mortal holes are then exactly
  the \emph{$\Pi$-holes} of one reader, its holes with an anchored comparison
  whose state pair lies in a relation $\Pi$. Mortality and the short rejected
  tail enter through the reader itself, and self events enter through its
  admission relation (\cref{sec:mo-anchored,sec:mo-separator}).
\item The short rejected tail bounds the depth of the separator word, by the
  depth theorem of \citet{Zhao26} (\cref{sec:mo-depth}).
\item We attach to each level of the hierarchy a \emph{positive channel},
  which records the anchored witness in the form in which it appears at that
  level, and show that channels have derivatives exactly as lettered monitors
  do (\cref{sec:mo-channels}). The bounded-depth search of that paper does not
  look at the channel and can lose the witness (\cref{ex:mo-loses}); once its
  exchange types, its shortening of the top level and its column summaries
  record the channel, it gives the anchored compression theorem
  (\cref{sec:mo-compression}).
\item Combining the three steps bounds the length of a shortest mortal hole by
  a doubly exponential function of the automaton state counts, hence a triply
  exponential function of $s$ for extracted instances, which decides the hole
  problem (\cref{sec:mo-short}).
\end{enumerate}
Condition (FH) is used only in the last theorem (\cref{rem:mo-fh}). The proofs
are in \cref{app:mo}.

\subsection{Anchored Comparisons}\label{sec:mo-anchored}

Fix a letter $\sep$, the \emph{separator}. Comparisons in this section are
comparisons of one word, in the end the separator word $u\sep t$ of
\cref{sec:mo-separator}, which is a different word from the pair $(u,t)$; we
therefore write them $(x,y)$ (here $x$ and $y$ are cuts, not variables),
whereas the cross comparisons $(s,e)$ of \cref{def:ex-events} relate cuts of
the two words of the pair. For a comparison $(x,y)$ of a word $w$, the factor $\wsub wx{x+\len w-y}=\wsuf wy$ is its \emph{matched factor},
and the letter $w[x+\len w-y]$ that follows the matched factor inside
$\wsuf wx$ is its \emph{continuation letter}; it exists because $x<y$. For a
terminal comparison $(x,\len w)$ the matched factor is empty and the
continuation letter is $w[x]$.

\begin{definition}[Anchored comparison, $\Pi$-hole]\label{def:mo-anchored}
A comparison of a word $w$ is \emph{anchored} if its continuation letter is
$\sep$. Let $(\mathcal M,\Adm,\Rej)$ be a reader with states $Q$ over a finite
alphabet $\Gamma\ni\sep$ (\cref{def:bg-reader}), and let $\Pi\subseteq Q\times
Q$. A \emph{$\Pi$-hole} of the reader is a nonempty hole $w$
(\cref{def:bg-hole}) that has an anchored comparison $(x,y)$ with
$(q_x,q_y)\in\Pi$.
\end{definition}

If $w$ contains $\sep$ exactly once, at the position $m$, then a pair of cuts
$(x,y)$ is an anchored comparison iff
\[
 x\le m,\qquad y=x+\len w-m,\qquad \wsuf wy=\wsub wxm .
\]
Indeed, the continuation letter of a comparison $(x,y)$ is $\sep$ iff
$x+\len w-y=m$, and then the matched factor $\wsuf wy$ has length $m-x$, so the
prefix relation $\wsuf wy\pre\wsuf wx$ says $\wsuf wy=\wsub wxm$; conversely
these three conditions make $(x,y)$ a comparison, since $x<y$ and $\wsub wxm$
is a prefix of $\wsuf wx$. The terminal pair $(m,\len w)$ is always anchored,
with empty matched factor. Every other anchored comparison matches a nonempty
factor that ends immediately before the separator with a suffix of the whole
word.

The equation $y=x+\len w-m$ is what matters. A comparison $(x,y)$ with $y>m$
only says that the suffix after $y$ is a \emph{prefix} of the suffix after
$x$; the equation forces the two matched factors to have the same length, and
the prefix relation becomes an equality. In the separator word this is exactly
the difference between an event, a prefix relation $\wsuf te\pre\wsuf ux$, and
a head event, an equality $\wsuf te=\wsuf ux$ (\cref{lem:mo-sep}). A
compression that preserves only ``some comparison with a state pair in $\Pi$''
is therefore not enough.

\begin{example}[An unanchored positive comparison]\label{ex:mo-unanchored}
Over $\Gamma=\{0,1,\sep\}$ take the reader whose states $0,\ldots,4$ form the
prefix chain of $01\sep0$, that is, $\delta(0,0)=1$, $\delta(1,1)=2$,
$\delta(2,\sep)=3$ and $\delta(3,0)=4$, with all other transitions into a
sink $5$. Let $\Rej=\{4\}$, $\Adm=\varnothing$ and $\Pi=\{(0,3)\}$. The only
rejected word is $w=01\sep0$, so it is the only hole. It has the comparison
$(0,3)$, since $\wsuf w3=0\pre w$, with state pair $(0,3)\in\Pi$, but the
continuation letter of $(0,3)$ is $w[1]=1$. The only anchored comparison of
$w$ is the terminal pair $(2,4)$, whose state pair $(2,4)$ is not in $\Pi$.
Hence the reader has no $\Pi$-hole, although its only hole has a comparison
labeled by $\Pi$.
\end{example}

\subsection{The Separator Word}\label{sec:mo-separator}

For words $u,t\in\Sigma^*$ put $w=u\sep t$, a word over
$\Gamma=\Sigma\sqcup\{\sep\}$, and $m=\len u$, so that $\len w=\len u+\len t+1$.
The cuts $x\le m$ of $w$ are the cuts of $u$, and the cut $m+1+e$ of $w$ is the
cut $e$ of $t$.

\begin{lemma}[Comparisons of the separator word]\label{lem:mo-sep}
Let $u,t\in\Sigma^*$, $w=u\sep t$ and $m=\len u$.
\begin{enumerate}[label=(\roman*),leftmargin=*]
\item No pair $(x,y)$ with $y\le m$ is a comparison of $w$.
\item For $x\le m$ and $e\le\len t$, the pair $(x,m+1+e)$ is a comparison of
  $w$ iff $(x,e)$ is a cross comparison of $(u,t)$.
\item For $i<j\le\len t$, the pair $(m+1+i,m+1+j)$ is a comparison of $w$ iff
  $(i,j)$ is a comparison of $t$.
\item For $x\le m$ and $e\le\len t$, the pair $(x,m+1+e)$ is an anchored
  comparison of $w$ iff $(x,e)$ is an aligned cross comparison of $(u,t)$, that
  is, iff $\wsuf ux=\wsuf te$.
\end{enumerate}
\end{lemma}

So every comparison of $w$ either joins a cut of the spine to a cut of the
target, and is then a cross comparison, or joins two cuts of the target
(\cref{fig:mo-separator}). The separator keeps the two words apart: every
matched factor is a suffix of $t$, so none runs across the separator, and a
matched factor that starts in $u$ stops either inside $u$ or right before the
separator; the comparisons of the second kind are the anchored ones. The proof
is in \cref{app:mo-separator}. All conditions of a mortal hole, mortality
included, are therefore read by one reader, which we now construct.

\begin{figure}[tbp]
\centering
\begin{tikzpicture}[>=Stealth,font=\small,x=0.8cm,y=1cm,
    every node/.style={inner sep=1.5pt},
    cell/.style={draw,minimum width=0.8cm,minimum height=0.5cm,inner sep=0pt},
    lb/.style={font=\scriptsize,inner sep=1pt},
    side/.style={font=\footnotesize,anchor=west,align=left}]
  % the word w = u # t
  \foreach \l/\f [count=\i from 0] in {0/white,1/black!12,0/black!12,\sep/white,1/black!12,0/black!12}
    \node[cell,fill=\f] (c\i) at (\i+0.5,0) {$\l$};
  \node[cell,very thick] at (3.5,0) {$\sep$};
  \foreach \i in {0,...,6}
    \node[lb,anchor=south] at (\i,0.27) {$\i$};
  \node[lb,anchor=east] at (-0.15,0.42) {cuts of $w$};
  \foreach \i in {0,...,3}
    \node[lb,anchor=north] at (\i,-0.27) {$\i$};
  \foreach \e in {0,1,2}
    \node[lb,anchor=north] at (4+\e,-0.27) {$\e$};
  \node[lb,anchor=east] at (-0.15,-0.42) {cuts of $u$, of $t$};
  \draw[decorate,decoration={brace,amplitude=3pt,mirror}] (0.05,-0.62) -- (2.95,-0.62)
    node[midway,below=3pt,lb]{$u$};
  \draw[decorate,decoration={brace,amplitude=3pt,mirror}] (4.05,-0.62) -- (5.95,-0.62)
    node[midway,below=3pt,lb]{$t$};
  % the anchored comparison (1,4)
  \draw[very thick,->] (1,0.72) .. controls (1.6,1.5) and (3.4,1.5) .. (4,0.72);
  % the cross comparison (0,5) and the self comparison (4,6)
  \draw[densely dashed,->] (0,0.72) .. controls (0.8,2.1) and (4.2,2.1) .. (5,0.72);
  \draw[densely dotted,thick,->] (4,0.72) .. controls (4.4,1.2) and (5.6,1.2) .. (6,0.72);
  % legend
  \draw[very thick] (7.0,1.55) -- (7.6,1.55);
  \node[side] at (7.7,1.55) {$(1,4)$: anchored, the aligned cross\\ comparison $(1,0)$: $\wsuf u1=t=10$};
  \draw[densely dashed] (7.0,0.6) -- (7.6,0.6);
  \node[side] at (7.7,0.6) {$(0,5)$: the cross comparison $(0,1)$:\\ $\wsuf t1=0\pre u$; continuation letter $1$};
  \draw[densely dotted,thick] (7.0,-0.35) -- (7.6,-0.35);
  \node[side] at (7.7,-0.35) {$(4,6)$: the self comparison\\ $(0,2)$ of $t$};
\end{tikzpicture}
\caption{The separator word $w=u\sep t$ for $u=010$ and $t=10$, the word of
\cref{ex:mo-loses-tower}; each arrow points from the first cut of a
comparison to its second. The cut $x\le3$ of $w$ is the cut $x$ of $u$, and the
cut $4+e$ is the cut $e$ of $t$. Every comparison of $w$ joins a cut of $u$ to
a cut of $t$, as the cross comparison $(0,5)$ does, or two cuts of $t$, as
$(4,6)$ does (\cref{lem:mo-sep}). The comparison $(1,4)$ is anchored: its
matched factor $\wsub w13=10$ (shaded) ends right before the separator and
equals the suffix $\wsuf w4$. This is how a head event, an equality of two
suffixes, appears in $w$; a cross comparison that is not aligned, such as
$(0,5)$, has a continuation letter inside $u$. The anchored comparisons of $w$
are $(1,4)$, $(2,5)$ and $(3,6)$.}
\Description{The word 0 1 0 separator 1 0 in six cells, the separator cell
framed thickly, with the cuts 0 to 6 of the word written above and, below,
the cuts 0 to 3 of u and 0 to 2 of t, with braces marking u and t. The two
cells 1 0 before the separator and the last two cells 1 0 are shaded. A thick
arrow goes from the cut 1 to the cut 4, a dashed arrow from the cut 0 to the
cut 5, and a dotted arrow from the cut 4 to the cut 6. A legend explains them:
the thick arrow is anchored and is the aligned cross comparison 1, 0; the
dashed arrow is the cross comparison 0, 1 with continuation letter 1; the
dotted arrow is the self comparison 0, 2 of t.}
\label{fig:mo-separator}
\end{figure}


\paragraph{The reader $\RPi$}
Mortality of a safe state is decided by reachability in the safe graph of $A$
(\cref{def:nm-kinds}). The reader $\RPi$ over $\Gamma=\Sigma\sqcup\{\sep\}$
reads $u\sep t$ in three phases.
\begin{itemize}[leftmargin=*]
\item On the spine, its states $\sPre(q,q')$, with $q\in S_A\setminus\Bad$
  and $q'\in S_{A'}$, follow $A$ and $A'$.
\item On the target up to its first rejected cut, its states
  $\sAct(r,r')$, with $r\in S_B$ and $r'\in S_{B'}$, follow $B$ and $B'$.
\item After the first rejected cut, its states $\sDone(r',k)$, with
  $r'\in S_{B'}$ and $1\le k\le N_A$, follow $B'$ and count the letters read,
  capped at $N_A$.
\end{itemize}
A state $\sSink$ absorbs every failure. The automata $A$, $B$, $A'$, $B'$
start in their state $0$ (\cref{def:ex-ss2self}); the start state of $\RPi$ is
$\sPre(0,0)$, or $\sSink$ if $0\in\Bad$. For a letter $i\in\Sigma$ the
transitions are
\[
\begin{aligned}
 \sPre(q,q')&\xrightarrow{i}\sPre(A(q,i),A'(q',i))
   \ \text{if }A(q,i)\notin\Bad,\text{ else }\sSink,\\
 \sPre(q,q')&\xrightarrow{\sep}\sAct(0,0)
   \ \text{if $q$ is mortal},\text{ else }\sSink,\\
 \sAct(r,r')&\xrightarrow{i}\sDone(B'(r',i),1)
   \ \text{if }r\in\Rej,\text{ else }\sAct(B(r,i),B'(r',i)),\\
 \sDone(r',k)&\xrightarrow{i}\sDone(B'(r',i),\min(k+1,N_A)),
\end{aligned}
\]
and every other transition, in particular every transition on $\sep$ after the separator,
leads to $\sSink$. The rejected states are $\sAct(r,r')$ with $r\in\Rej$ and
$\sDone(r',k)$ with $k\le N_A-1$. The admission relation consists of the pairs
\[
 \bigl(\sPre(q,q'),\sAct(r,r')\bigr)\ \text{with }(q,r)\in\Guard,
 \qquad
 \bigl(\sAct(r_1,r'_1),\sAct(r_2,r'_2)\bigr)\ \text{with }(r_1,r_2)\in\GuardS,
\]
and no others. The relation $\Pi$ relates $\sPre(q,q')$ to $\sAct(r,r')$ and to
$\sDone(r',k)$ iff $(q',r')\in\Head$, and relates no other pairs.

Thus $\RPi$ reads the separator only from a mortal state, rejects a word only if
fewer than $N_A$ letters follow the first rejected cut of the target, and lets
$\Pi$ compare the head automata on the two sides of the separator. The target
automaton is kept up to and including the first rejected cut, because events
and self events may end there; after it, events and self events no longer
exist (\cref{def:ex-events}), so only $B'$ and the counter remain. The admissions
from the self guard $\GuardS$ are the only difference from a reader for
instances without self events; the compression below allows an arbitrary admission relation.

For a reader, only the states reachable from its start state matter: the
restriction to them is again a reader, with the same run on every word and
hence the same holes and $\Pi$-holes. Let $r_I$ and $z_I$ be the numbers of
states of the product automata $A\times A'$ and $B\times B'$ reachable from
their start states, so that $r_I\le N_AN_{A'}$ and $z_I\le N_BN_{B'}$, and put
\[
 \Nbar=1+r_I+z_I+N_AN_{B'},\qquad d_I=2\Nbar N_A-1 .
\]

\begin{lemma}[The $\Pi$-reader]\label{lem:mo-pi-reader}
The reader $\RPi$ and the relation $\Pi$ are computable from $I$, and:
\begin{enumerate}[label=(\roman*),leftmargin=*]
\item at most $\Nbar$ states of $\RPi$ are reachable from its start state;
\item a word $w$ over $\Sigma\sqcup\{\sep\}$ is a $\Pi$-hole of $\RPi$ iff
  $w=u\sep t$ with $u,t\in\Sigma^*$ and $(u,t)$ a mortal hole of $I$;
\item for such $w$ the first rejected cut is $\len u+1+J(t)$, so the rejected
  tail of $w$ is $\len t-J(t)$.
\end{enumerate}
\end{lemma}

On $u\sep t$, the state at a cut $x\le\len u$ carries $A_x(u)$ and $A'_x(u)$
while $u$ is safe, and the state at the cut $\len u+1+e$ carries $B_e(t)$ up to
$J(t)$ and $B'_e(t)$ throughout. By \cref{lem:mo-sep}, the admitted comparisons
are then exactly the events and the self events, and $\Pi$ on the anchored
comparisons is exactly the head condition. The proof is in
\cref{app:mo-separator}.

\subsection{The Depth of the Separator Word}\label{sec:mo-depth}

The length bound of the bounded-depth search grows with the depth of the word
(\cref{thm:bg-search}). For a mortal hole, the depth of the separator word is
bounded at once.

\begin{lemma}[Depth of mortal holes]\label{lem:mo-depth}
Every mortal hole $(u,t)$ of $I$ satisfies $\depth(u\sep t)\le d_I$.
\end{lemma}

\begin{proof}
By \cref{lem:mo-pi-reader}, $u\sep t$ is a nonempty rejected word of the
reachable part of $\RPi$, a reader with at most $\Nbar$ states, and its rejected
tail is $\len t-J(t)\le N_A-1$. \Cref{thm:bg-depth} gives
$\depth(u\sep t)\le2\Nbar((N_A-1)+1)-1=d_I$.
\end{proof}

The bound is coarse, since it involves the spine and head automata through
$\Nbar$, but it enters the final estimate only through a polynomial in $\Nbar$
and $N_A$, and changes only the constant in the exponent
(\cref{lem:mo-estimate}).

\subsection{Positive Channels}\label{sec:mo-channels}

The bounded-depth search sees each level of the hierarchy through a lettered
monitor, the derivative of the monitor one level below
(\cref{def:bg-monitor,def:bg-derivative}); the derivatives carry the negative
conditions, because a comparison of one level reappears at the next as a
comparison of blocks or as a test in the final target. We carry the positive
condition in the same way. At level $0$ it is an anchored comparison in $\Pi$,
a comparison together with a test of its continuation letter; one level up, the
continuation letter is the first letter of the expansion of a block, or a
letter inside the last block, and the comparison may have become a test of the
final label. Five components cover all the forms that arise.

\begin{definition}[Positive channel]\label{def:mo-channel}
Let $\mathcal M$ be a lettered monitor over an alphabet $\Gamma$ with state set
$\mathcal Q$. A \emph{positive channel} $\Theta=(F,R,S,S',\Phi)$ on $\mathcal M$
consists of a predicate $F\subseteq\mathcal Q\times\Gamma$, three relations
$R$, $S$ and $S'$ on labels, and a set $\Phi\subseteq\Gamma$ of letters. For a
nonempty word $v$ with labels $L_0,\ldots,L_{\len v}$, the \emph{positive
condition} $\PosOp_\Theta(v)$ holds iff one of the following clauses holds:
\begin{description}[leftmargin=2.8em,font=\normalfont]
\item[(F)] $F(L_{\len v})$;
\item[(R)] some comparison $(x,y)$ of $v$ with $y<\len v$ has $R(L_x,L_y)$ and
  $v[x+\len v-y]\in\Phi$;
\item[(S)] some cut $x<\len v$ has $S(L_x,L_{\len v})$;
\item[(S$'$)] some cut $x<\len v$ has $S'(L_x,L_{\len v})$ and $v[x]\in\Phi$.
\end{description}
A \emph{positive hole} of $(\mathcal M,\Theta)$ is a hole $v$ of $\mathcal M$
with $\PosOp_\Theta(v)$.
\end{definition}

Clauses (R) and (S$'$) together say that some comparison has its labels in a
relation, $R$ if it is not terminal and $S'$ if it is, and its continuation
letter in $\Phi$; the continuation letter of the terminal comparison $(x,\len
v)$ is $v[x]$. Clauses (F) and (S) are tests without a continuation letter.
They are empty at level $0$, but derivatives move witnesses into them, as the
derivatives of lettered monitors move admissions into the terminal relation and the
final target.

For a reader $(\mathcal M,\Adm,\Rej)$ over $\Gamma\ni\sep$ with start $q_0$ and
$\Pi\subseteq Q\times Q$, the \emph{base channel} on the lettered monitor of
the reader (\cref{def:bg-monitor}) is
\[
 \Theta_\Pi=(\varnothing,\widehat\Pi,\varnothing,\widehat\Pi,\{\sep\}),
 \qquad
 \widehat\Pi(L,L')\iff(\den L,\den L')\in\Pi .
\]

\begin{lemma}[The base channel]\label{lem:mo-base}
A nonempty word $w$ satisfies $\PosOp_{\Theta_\Pi}(w)$ iff it has an anchored
comparison $(x,y)$ with $(q_x,q_y)\in\Pi$. Hence the $\Pi$-holes of the reader
are exactly the positive holes of its lettered monitor with the base channel.
\end{lemma}

\begin{proof}
Since $\den L_i=q_i$, clause (R) is an anchored comparison $(x,y)$ with
$y<\len w$ and $(q_x,q_y)\in\Pi$, and clause (S$'$) is the same for a terminal
comparison $(x,\len w)$, whose continuation letter is $w[x]$; clauses (F) and
(S) are empty. The reader and its lettered monitor have the same nonempty
holes.
\end{proof}

The derivative of a channel uses the shifted blocks of
\cref{def:bg-derivative}. The letter that follows the position $g$ of the
shifted block of $X$ is $(X\ell)[g]$: the letter $X[g]$ for $g<\len X$, and
the marker that opens the next shifted block for $g=\len X$.

\begin{definition}[Derivative of a channel]\label{def:mo-channel-derivative}
Let $\Theta=(F,R,S,S',\Phi)$ be a positive channel on $\mathcal M$ and
$\ell\in\Gamma$. The \emph{derivative} $\Theta_\ell=(F_\ell,R_\ell,S_\ell,
S'_\ell,\Phi_\ell)$ is the positive channel on $\mathcal M_\ell$ given by
\[
\begin{aligned}
 \Phi_\ell&=\{X:(X\ell)[0]\in\Phi\},\\
 R_\ell((\zeta,X),(\zeta',X'))&\iff X=xy,\ X'=x'y,\
   R\bigl(\lab_{\len x}(\zeta,X),\lab_{\len{x'}}(\zeta',X')\bigr)
   \text{ for some }x,x',y,\\
 S'_\ell&=R_\ell,\\
 S_\ell((\zeta,X),(\zeta',X'))&\iff\ptest(\zeta,X;\delta_\ell(\zeta',X')),\\
 F_\ell(\zeta,X)&\iff F(\delta_\ell(\zeta,X),\ell)\ \text{or}\
   \ptest(\zeta,X;\delta_\ell(\zeta,X)),
\end{aligned}
\]
where the \emph{positive terminal test} is
\[
\begin{aligned}
 \ptest(\zeta,X;\zeta'')\iff{}&\text{for some }g\le\len X:\
 S\bigl(\lab_g(\zeta,X),(\zeta'',\ell)\bigr),\\
 &\text{or }S'\bigl(\lab_g(\zeta,X),(\zeta'',\ell)\bigr)\text{ and }
 (X\ell)[g]\in\Phi,
\end{aligned}
\]
and every relation with the start label $\triangleright$ of $\mathcal M_\ell$ as
an argument is false.
\end{definition}

The relation $R_\ell$ is built from $R$ exactly as the derived admission
$\Lambda_\ell$ is built from $\Lambda$, and the positive terminal test is the
analog of the terminal test in $\Lfin_\ell$ and $\Tfin_\ell$: it relates a
position of a block, possibly the closing marker, to the final label, which is
$(\zeta'',\ell)$ when the final marker is read from the state $\zeta''$. The predicate $F$ concerns
only the final label, so it enters $F_\ell$ but not $S_\ell$. The set
$\Phi_\ell$ tests the first letter of the expansion of a block, because the
continuation letter of a comparison of the expansion is the first letter of the
expansion of the continuation letter of the corresponding comparison of blocks
(\cref{lem:mo-transport}). This is how the anchoring equation travels up the
hierarchy: ``followed by $\sep$'' at level $0$ becomes ``followed by a block
whose expansion starts with $\sep$'' at level $1$, and so on.

\begin{lemma}[Positive desubstitution]\label{lem:mo-desub}
Let $\Theta$ be a positive channel on $\mathcal M$, let $\ell\in\Gamma$, and let
$v$ be the expansion of a nonempty word $v'$ of blocks avoiding $\ell$. Then
$\PosOp_\Theta(v)\iff\PosOp_{\Theta_\ell}(v')$. Consequently $v$ is a positive
hole of $(\mathcal M,\Theta)$ iff $v'$ is a positive hole of
$(\mathcal M_\ell,\Theta_\ell)$.
\end{lemma}

The second claim follows from the first and \cref{lem:bg-desub}. The proof of
the first, in \cref{app:mo-desub}, translates each clause by the label facts
\cref{eq:bg-label-facts} and the alignment lemma (\cref{lem:bg-alignment}): a
nonterminal comparison of $v$ becomes a comparison of blocks, read by $R_\ell$
or, if it ends at the last cut, by $S'_\ell$; a terminal comparison of $v$ that
starts in the block $X_i$ becomes a test of $S_\ell$ or, if $X_i$ is the last
block, of $F_\ell$.

\begin{corollary}[Channel towers]\label{cor:mo-tower}
Let $\mathcal M_0$ be a lettered monitor with $N$ states, $\Theta_0$ a positive
channel on it, and $w$ a nonempty word with canonical hierarchy
$(w_j,\ell_j)$. Put $\mathcal M_{j+1}=(\mathcal M_j)_{\ell_j}$ and
$\Theta_{j+1}=(\Theta_j)_{\ell_j}$. Then $w$ is a positive hole of
$(\mathcal M_0,\Theta_0)$ iff $w_j$ is a positive hole of
$(\mathcal M_j,\Theta_j)$, for every $j$.
\end{corollary}

\begin{proof}
Each $w_{j+1}$ is a nonempty word of blocks avoiding $\ell_j$, with expansion
$w_j$. Apply \cref{lem:mo-desub} at each level.
\end{proof}

Neither the lemma nor the corollary assumes that $\sep$ occurs once, or at
all, and the reader may have an arbitrary admission relation.

\begin{example}[A witness travels up the tower]\label{ex:mo-loses-tower}
Take $\Gamma=\{0,1,\sep\}$, the five states $0,\ldots,4$ with start $0$, the
transitions
\begin{center}
\begin{tabular}{@{}c|ccccc@{}}
state & $0$ & $1$ & $2$ & $3$ & $4$\\
\hline
on $0$ & $3$ & $0$ & $0$ & $3$ & $0$\\
on $1$ & $1$ & $1$ & $4$ & $0$ & $2$\\
on $\sep$ & $4$ & $3$ & $2$ & $4$ & $0$
\end{tabular}
\end{center}
$\Adm=\varnothing$, $\Rej=\{0,1\}$ and $\Pi=\{(3,4),(4,2)\}$. The word
$w=010\sep10$ has the run $0,3,0,3,4,2,0$ and ends in a rejected state, so it is
a hole. Its anchored comparisons are $(1,4)$, $(2,5)$ and $(3,6)$, with state
pairs $(3,4)$, $(0,2)$ and $(3,0)$. Only the first, which matches the factor
$10$ before the separator with the final $10$, is in $\Pi$, so $w$ is a
$\Pi$-hole; for the base channel, its witness $(1,4)$ is clause (R). The
hierarchy of $w$ is $w_1=abc$ with $a=[\eps]$, $b=[1]$ and $c=[\sep1]$, and
$w_2=[ab]$, a single letter. At level $1$, with the marker $0$, the cut $1$
lies in the block $X_1=1$ and the cut $4$ in the last block $X_2=\sep1$, both
with the common suffix $1$; so the comparison of blocks is the terminal pair
$(2,3)$ of $w_1$, and the witness becomes clause (S$'$) at the cut $2$, with
continuation letter $c$, which lies in $\Phi_1$ because its expansion starts
with $\sep$. At level $2$, with the marker $c$, the cut $2$ of $w_1$ is the end
of the only block $ab$ of $w_2$; so the witness becomes the second disjunct of
the final target $F_2$, through the positive terminal test at the position
$\len{ab}$, whose letter is the marker $c\in\Phi_1$. Clause (R) at level $0$
has become clause (S$'$) at level $1$ and clause (F) at level $2$.
\end{example}

\Cref{fig:mo-tower} draws the three levels of this example. Each level sees
the witness only through its own letters: once the matched factor $10$ has
become the single letter $b$ and the separator the first letter of $c$, the
comparison turns into a test at the boundary between $b$ and $c$, and one
level higher into a test of the final label.

\begin{figure}[tbp]
\centering
\begin{tikzpicture}[>=Stealth,font=\small,x=0.75cm,y=1cm,
    every node/.style={inner sep=1.5pt},
    cell/.style={draw,minimum height=0.5cm,inner sep=0pt},
    lb/.style={font=\scriptsize,inner sep=1pt},
    side/.style={font=\footnotesize,anchor=west,align=left},
    lvl/.style={font=\footnotesize,anchor=east}]
  % level 0: w = 010#10 with its run
  \foreach \l/\f [count=\i from 0] in {0/white,1/black!12,0/black!12,\sep/white,1/black!12,0/black!12}
    \draw[fill=\f] (\i,-0.25) rectangle node{$\l$} (\i+1,0.25);
  \foreach \q [count=\i from 0] in {0,3,0,3,4,2,0}
    \node[lb,anchor=north] at (\i,-0.28) {$\q$};
  \draw[very thick,->] (1,-0.6) .. controls (1.6,-1.25) and (3.4,-1.25) .. (4,-0.6);
  \node[lvl] at (-0.3,0) {level $0$: $w$};
  \node[lvl,font=\scriptsize] at (-0.3,-0.42) {run};
  \node[side] at (6.5,-0.45) {(R): the anchored comparison $(1,4)$,\\
    state pair $(3,4)\in\Pi$, continuation letter $\sep$};
  % level 1: w_1 = abc, marker 0
  \draw (0,1.05) rectangle node{$a$} (1,1.55);
  \draw (1,1.05) rectangle node{$b=[1]$} (3,1.55);
  \draw (3,1.05) rectangle node{$c=[\sep1]$} (6,1.55);
  \node[lb,anchor=south] at (0.5,1.55) {$[\eps]$};
  \foreach \x in {0,1,3,6} \draw[densely dotted] (\x,0.25) -- (\x,1.05);
  \draw[very thick] (3,1.05) -- (3,1.55);
  \node[lvl] at (-0.3,1.3) {level $1$: $w_1$};
  \node[side] at (6.5,1.3) {(S$'$) at the cut $2$ of $w_1$, the terminal comparison\\
    $(2,3)$, continuation letter $c\in\Phi_1$};
  % level 2: w_2 = [ab], marker c
  \draw (0,2.1) rectangle node{$[ab]$} (6,2.6);
  \foreach \x in {0,6} \draw[densely dotted] (\x,1.55) -- (\x,2.1);
  \node[lvl] at (-0.3,2.35) {level $2$: $w_2$};
  \node[side] at (6.5,2.35) {(F): the positive terminal test at the\\
    position $\len{ab}$ of the block, letter $c\in\Phi_1$};
  % after deleting the column of b
  \begin{scope}[yshift=-2.35cm]
    \foreach \l/\f [count=\i from 0] in {0/white,\sep/white,1/white,0/white}
      \draw[fill=\f] (\i,-0.25) rectangle node{$\l$} (\i+1,0.25);
    \foreach \q [count=\i from 0] in {0,3,4,2,0}
      \node[lb,anchor=north] at (\i,-0.28) {$\q$};
    \draw[densely dashed,->] (0,0.3) .. controls (0.6,0.85) and (2.4,0.85) .. (3,0.3);
    \draw[densely dashed,->] (1,0.3) .. controls (1.6,0.85) and (3.4,0.85) .. (4,0.3);
    \node[lvl] at (-0.3,0) {without $b$};
    \node[side] at (6.5,0.05) {anchored comparisons $(0,3)$ and $(1,4)$,\\
      state pairs $(0,2)$ and $(3,0)$, both outside $\Pi$};
  \end{scope}
\end{tikzpicture}
\caption{A positive witness travels up the hierarchy
(\cref{ex:mo-loses-tower,ex:mo-loses}). Bottom up: the hole $w=010\sep10$ of
the reader of \cref{ex:mo-loses-tower} with its run, the word $w_1=abc$ of
blocks for the marker $0$, each letter drawn over its block and the marker that
closes it, and the one-letter word $w_2=[ab]$ for the marker $c$. The anchored
witness $(1,4)$ in $\Pi$, whose matched factor $10$ is shaded, is clause (R) of
the base channel at level $0$, clause (S$'$) of the derived channel at level
$1$, and clause (F) at level $2$, where it is a test of the final label
(\cref{def:mo-channel,lem:mo-desub}). The compression for plain entailment sees
no difference between the prefixes $a$ and $ab$ of the block and deletes the
column of $b$, which carries the matched factor $10$; the result $0\sep10$
(last row) is still a hole, but no longer a $\Pi$-hole. Recording the channel in
the summaries keeps $b$.}
\Description{Three stacked rows of cells. The bottom row is the word 0 1 0
separator 1 0, with the states 0 3 0 3 4 2 0 of its run under the cuts and a
thick arrow below from the cut 1 to the cut 4; the cells 1 0 before the
separator and the last cells 1 0 are shaded. The middle row has the letters a
over the first cell, b equal to the block 1 over the next two cells, and c
equal to the block separator 1 over the last three cells; the boundary between
b and c, the cut 2, is drawn thick. The top row is one cell, the block a b,
over the whole word. Notes to the right name the clause of the witness at each
level: R at level 0, S prime at level 1, F at level 2. A separate row at the
bottom shows the word 0 separator 1 0 with the run 0 3 4 2 0 and two dashed
arrows, from the cut 0 to the cut 3 and from the cut 1 to the cut 4, whose
state pairs are outside Pi.}
\label{fig:mo-tower}
\end{figure}


\subsection{Anchored Compression}\label{sec:mo-compression}

\paragraph{What the compression for plain entailment misses}
Steps~2 and~4 of the bounded-depth search (\cref{sec:bg-hierarchy}, and
precisely \cref{app:bg-search}) replace the word by another one: step~2 merges
isolated letters with the same exchange type, and step~4 deletes the columns
between two cuts of a column word with equal summaries. Each of them only has
to keep the word a hole. A $\Pi$-hole must also \emph{contain} an anchored
comparison in $\Pi$. The exchange types and summaries record nothing about
$\Pi$, so a column deletion, and similarly a merge, can destroy every such
comparison.

\begin{example}[The compression for plain entailment loses the anchored witness]\label{ex:mo-loses}
Continue \cref{ex:mo-loses-tower}. Step~4 shortens the block $ab$ over the
level-$1$ monitor with the marker $c$. In this monitor $c$ leads every state to
$2$, and the exits of the shifted blocks of the prefixes $a$ and $ab$ are $0$
from every entry state; since $\Adm$ is empty, all terminal and admission bits
are false. So the summaries after $a$ and after $ab$ coincide, and step~4
deletes the column holding $b$. Expanding $[a]$ gives $0\sep10$
(\cref{fig:mo-tower}, last row), with run
$0,3,4,2,0$: a hole, whose anchored comparisons $(0,3)$ and $(1,4)$ have the
state pairs $(0,2)$ and $(3,0)$, both outside $\Pi$. The deleted letter $b$ is
exactly the letter whose expansion is the matched factor $10$ before the
separator.
\end{example}

The bound of the next theorem is an explicit function $\Lensharp(N,d)$. Like
$\holeBound(N,d)$ of \citet{Zhao26}, it is built from constants $\sCsharp$,
$\sAsharp$, $\sGsharp$ and $\scolsharp$, which bound the anchored exchange
types, the protected and isolated letters, the dictionaries and the positive
column summaries. They are defined in \cref{eq:mo-constants}, increase in each
argument, and do not depend on the alphabet; the theorem states their closed
form.

\begin{theorem}[Anchored compression]\label{thm:mo-anchored}
Let $(\mathcal M,\Adm,\Rej)$ be a reader with $N\ge1$ states over a finite
alphabet $\Gamma\ni\sep$, let $\Pi\subseteq Q\times Q$, and let $d\ge0$. If the
reader has a $\Pi$-hole of depth at most $d$, then it has a $\Pi$-hole of depth
at most $d$ and length at most $\Lensharp(N,d)$. More generally, for every
lettered monitor $\mathcal M_0$ with $N$ states over any alphabet and every
positive channel $\Theta_0$ on it, if $(\mathcal M_0,\Theta_0)$ has a positive
hole of depth at most $d$, then it has one of depth at most $d$ and length at
most $\Lensharp(N,d)$. Moreover, with
$\Xisharp(N,d)=6N^2+10N(N+d)(d+1)+7N+3d+4$,
\[
 \sGsharp(N,d)\le2^{\Xisharp(N,d)}
 \qquad\text{and}\qquad
 \Lensharp(N,d)\le2^{2^{5\Xisharp(N,d)}} .
\]
\end{theorem}

The theorem is the analog of \cref{thm:bg-search} for $\Pi$-holes. It
preserves the \emph{existence} of an anchored comparison in $\Pi$, not its
coordinates or its state pair; to keep a particular state pair, apply it to
the singleton relation of that pair. It needs no uniqueness of the separator.
The first claim is the second one for the lettered monitor of the reader and
the base channel, by \cref{lem:mo-base}.

\paragraph{Why recording the channel suffices}
The positive condition is a property of the labels at the cuts of a witness
and of one letter, the continuation letter. The shortenings of \citet{Zhao26}
(steps~2--4 in \cref{sec:bg-hierarchy}) preserve such data once they record
the channel. The merges preserve the state run, the comparison positions, and
the relevant channel tests; the top shortening preserves the labels of the
cuts of the witness, which it never deletes, and the letter, which is constant
on a unary word; and the
summaries determine the derived channel, so that agreement preserves labels
and letters up to renaming. This is why recording the channel suffices, with
no further condition on the separator.

\begin{proof}[Proof outline]
The complete proof is in \cref{app:mo-proof}. The proof of \cref{thm:bg-search} merges letters
with equal exchange types and deletes columns between equal summaries. We add
the five components of the channel to both, so that each shortening keeps some
anchored witness, though not necessarily the original one; every argument
about the negative conditions is reused unchanged. Let $w$ be a positive hole
of depth $D\le d$, with hierarchy $(w_j,\ell_j)$; we follow the five steps of
that proof (\cref{sec:bg-hierarchy}).
\begin{enumerate}[leftmargin=*]
\item \emph{Towers.} By \cref{cor:mo-tower}, $w_j$ is a positive hole of
  $(\mathcal M_j,\Theta_j)$ for every $j$; the clause of the witness may change
  from level to level (\cref{ex:mo-loses-tower}).
\item \emph{Merging letters.} In the terms of step~2 in \cref{sec:bg-hierarchy},
  a letter of $w_i$ is unprotected if it is neither the marker nor a letter of
  the start prefix, and isolated if, moreover, all its unequal comparison
  partners are protected; its exchange type is the table of how $\mathcal M_i$
  treats it when it is read from an ordinary state. The
  \emph{anchored exchange type} of an unprotected letter $a$ at a level
  $1\le i<D$ adds the same table for the channel: the row of $F_i$ at $a$, the
  self table, the protected rows and the initial row of each of $R_i$, $S_i$
  and $S'_i$, and the bit $[a\in\Phi_i]$ (\cref{def:mo-type}). There are at
  most $\sCsharp(N,d)$ anchored types (\cref{lem:mo-types}). Merging two
  isolated letters with the same anchored type keeps the comparisons and the
  run of the level, so it keeps every value of $\Lambda_i$, $\Lfin_i$, $R_i$,
  $S_i$ and $S'_i$ at the comparisons of the level, exactly as in the
  argument of \citet{Zhao26} for $\Lambda_i$ and $\Lfin_i$; it keeps the final label,
  because the marker is protected; and it keeps the position of every
  continuation letter and changes that letter at most from $a$ to $b$, which
  have the same bit. So the merged level is again positive, and expanding it
  gives a positive hole of $(\mathcal M_0,\Theta_0)$ of the same depth, no
  longer, with a smaller dictionary (\cref{lem:mo-exchange}). Repeated merges
  bound the dictionary, the set of letters, of every level $i\ge1$ by
  $\sGsharp(N,d)$ (\cref{lem:mo-dictionary}). A merge is not injective and
  gives no agreement; after each merge the hierarchy and the towers are
  recomputed.
\item \emph{The top level.} The top level $w_D$ is unary. Deleting a loop
  between two cuts with equal states keeps a hole (\cref{lem:bg-unary}), and it
  keeps the positive witness if no cut of the witness is deleted. Protecting the
  at most three cuts of one witness leaves at most three stretches of
  pairwise distinct states, so $\len{w_D}\le3(N+D)$ after shortening
  (\cref{lem:mo-unary}).
\item \emph{Deleting columns.} The \emph{positive summary} of a tuple of blocks
  extends the summary of \cref{app:bg-search} by fields that record the positive terminal
  test, the last pending $S'$-test (a test at the end of a block, whose letter
  is the next one to be read), whether a block is empty, the bit $\Phi$ of its
  first letter, and the lift $R_\ell$ between pairs of blocks. These fields are
  updated column by column and determine the derived channel on the blocks, in
  particular $S_\ell$ without an $F$-term and $F_\ell$ with one
  (\cref{lem:mo-summary}). Deleting the columns between two cuts with equal
  positive summaries therefore keeps both the derived monitor and the derived
  channel (\cref{lem:mo-columns}). In \cref{ex:mo-loses}, the pending
  $S'$-test is false after $a$ and true after $ab$, where it records the
  witness $(1,4)$; so the positive summaries after $a$ and after $ab$ differ,
  and the column holding $b$ is kept (\cref{ex:mo-loses-fields}).
\item \emph{Agreement and the bound.} Agreement of monitors on the letters in
  use (step~4 in \cref{sec:bg-hierarchy}) is extended to the five components
  of the channel, which must also take equal values at corresponding letters
  and labels. Extended agreement preserves positive holes and is preserved
  under derivation (\cref{lem:mo-agreement}). The supported compression of
  \citet{Zhao26}, which shortens the letters in use level by level from the top,
  then runs unchanged with the positive summaries (\cref{lem:mo-compression}).
  Each of the at most $d$ expansions costs a factor
  $\scolsharp(N+d+1,\sGsharp(N,d))$, and the top level costs $3(N+d+1)$, which
  gives $\Lensharp(N,d)$. The new top word is a positive hole of the new top
  monitor, and \cref{lem:bg-desub,lem:mo-desub} reconstruct a positive hole of
  $(\mathcal M_0,\Theta_0)$ of depth at most $D$.
\end{enumerate}
The closed-form estimate is \cref{lem:mo-closed-form}.
\end{proof}

\subsection{Short Mortal Holes and Decidability}\label{sec:mo-short}

Put
\[
 \Lensharp_I=\Lensharp(\Nbar,d_I),
\]
with $\Nbar$ and $d_I$ as in \cref{sec:mo-separator}, and recall the bound
$\Lnm$ of \cref{thm:nm-three-cases}, evaluated at the state counts
$(N_A,N_B,N_{A'},N_{B'})$.

\begin{theorem}[Short mortal holes]\label{thm:mo-short}
If $I$ has a mortal hole, then it has a mortal hole $(u,t)$ with
$\len u+1+\len t\le\Lensharp_I$. No assumption on $I$ is needed; in particular
(FH) is not used.
\end{theorem}

\begin{proof}
Let $(u,t)$ be a mortal hole. By \cref{lem:mo-pi-reader}, $w=u\sep t$ is a
$\Pi$-hole of $\RPi$, and of its reachable part, a reader with at most $\Nbar$
states over the finite alphabet $\Sigma\sqcup\{\sep\}$. By \cref{lem:mo-depth},
$\depth(w)\le d_I$. By \cref{thm:mo-anchored} and monotonicity of $\Lensharp$,
the reachable part has a $\Pi$-hole $\tilde w$ of length at most
$\Lensharp_I$. By \cref{lem:mo-pi-reader}(ii), $\tilde w=\tilde u\sep\tilde t$
with $(\tilde u,\tilde t)$ a mortal hole, and
$\len{\tilde u}+1+\len{\tilde t}=\len{\tilde w}$.
\end{proof}

\begin{corollary}[Mortal holes are decidable]\label{cor:mo-decidable}
It is decidable whether an instance of the two-word hole problem with self
events has a mortal hole.
\end{corollary}

\begin{theorem}[Short holes]\label{thm:mo-decide}
Let $I$ satisfy (FH). Then $I$ has a hole iff it has a hole $(u,t)$ with
\[
 \len u+\len t<\max\bigl(\Lnm(N_A,N_B,N_{A'},N_{B'}),\ \Lensharp_I\bigr).
\]
In particular, it is decidable whether an instance with forced target heads
has a hole.
\end{theorem}

\begin{proof}
By \cref{cor:nm-short-holes}, $I$ has a hole iff it has a hole with
$\len u+\len t<\Lnm(N_A,N_B,N_{A'},N_{B'})$ or a mortal hole. In the second case
\cref{thm:mo-short} gives a mortal hole, which is a hole, with
$\len u+\len t<\Lensharp_I$. The bound is computable from $I$, and the alphabet
$\Sigma$ is finite, so there are finitely many candidate pairs; for each, the
conditions of \cref{def:ex-ss2self} are checked over finitely many cuts.
\end{proof}

The proof of \cref{cor:mo-decidable} is the same enumeration with the bound of
\cref{thm:mo-short} and the additional tests that $A(u)$ is mortal and
$\len t-J(t)<N_A$ (\cref{app:mo-estimates}). Both bounds are independent of the
alphabet, and hence of the arity; the alphabet enters only the number of
candidate pairs. For $n=0$ the only candidate is $(\eps,\eps)$.

\paragraph{The size of the bound}
Write $E_1(x)=2^x$ and $E_{k+1}(x)=2^{E_k(x)}$.

\begin{lemma}[Mortal bound]\label{lem:mo-estimate}
Let $\alpha\ge3$ and $\beta\ge0$ be integers with $\Nbar\le2^\alpha$ and
$N_A\le2^\beta$. Then $\Lensharp_I\le E_3(3\alpha+2\beta+9)$.
\end{lemma}

\begin{corollary}[Mortal bound for extracted instances]\label{cor:mo-extracted}
Let $\varphi$ be a flat system over $X$ and $\psi$ a flat system over $X\cup Z$,
with $v=\len X$, $h=\len Z$ and $s=v+h$. Every extracted instance $I_{pq}$
(\cref{def:ex-instance}) satisfies
\[
 \Lensharp_{I_{pq}}\le E_3(10v+6s+18)\le E_3(16s+18),
\]
for every arity $n$ and every variance. So if $I_{pq}$ has a mortal hole, it
has one whose separator word has length at most $E_3(16s+18)$.
\end{corollary}

\begin{proof}
By \cref{sec:ex-instances}, $N_A=2^{2v}$, $N_B=2^{4v+2h+1}=2^{2v+2s+1}$ and
$N_{A'}=N_{B'}=2^{2s}$, and the reachable pairs of states of $B$ and $B'$ number
at most $N_B$. Hence $r_I\le N_AN_{A'}=2^{2v+2s}$, $z_I\le2^{2v+2s+1}$ and
$N_AN_{B'}=2^{2v+2s}$, so $\Nbar\le1+2^{2v+2s+2}\le2^{2v+2s+3}$.
\Cref{lem:mo-estimate} with $\alpha=2v+2s+3$ and $\beta=2v$ gives
$E_3(3(2v+2s+3)+4v+9)=E_3(10v+6s+18)$, and $v\le s$. The second claim is
\cref{thm:mo-short}. None of these sizes depends on $n$.
\end{proof}

\begin{remark}[Where (FH) is used]\label{rem:mo-fh}
Condition (FH) is used only in \cref{thm:mo-decide}, through
\cref{thm:nm-three-cases}. Every other result of this section, in particular
\cref{thm:mo-short,cor:mo-decidable}, holds for every instance.
\end{remark}
